% Source problem: https://www.ap.t.u-tokyo.ac.jp/themes/ap_t_u-tokyo/graduate/pdf/h31ap_exam.pdf
\documentclass[11pt,a4paper]{article}
\usepackage[margin=24mm]{geometry}
\usepackage{fontspec}
\usepackage{xeCJK}
\setCJKmainfont{HaranoAjiMincho-Regular.otf}
\usepackage{amsmath,amssymb}
\setlength{\parindent}{1em}
\setlength{\parskip}{0.35em}
\title{物理工学専攻 院試解答集\\2019年度 物理学I 第2問}
\author{政岡凜太郎}
\date{}
\begin{document}
\maketitle

\section*{[1.1] 誘導電場}
半径$r$の円筒経路を上から見て反時計回りを正に取る。Faradayの法則から
\begin{equation}
 \oint\boldsymbol E\cdot d\boldsymbol l=-\frac{d}{dt}(\pi r^2B(t)),\qquad
 2\pi rE_\theta=-\pi r^2\dot B(t).
\end{equation}
$B(t)=\mu_0H_0\sin(2\pi ft)$より
\begin{equation}
 \boxed{E_\theta(r,t)=-\frac r2\dot B(t)
 =-\pi f\mu_0H_0r\cos(2\pi ft)}.
\end{equation}

\section*{[1.2] 誘導加熱電力}
電流密度は$J=E_\theta/\rho$なので、単位体積あたりのJoule発熱率は$E_\theta^2/\rho$である。体積要素$2\pi rL\,dr$にわたって積分すると
\begin{align}
 P(t)&=\int_0^R\frac{E_\theta(r,t)^2}{\rho}2\pi rL\,dr\\
 &=\frac{\pi^3\mu_0^2H_0^2f^2LR^4}{2\rho}\cos^2(2\pi ft).
\end{align}

\section*{[1.3] 一秒間の温度上昇}
$1$秒は$f=10\,\mathrm{Hz}$の整数周期なので、$\cos^2$の時間平均は$1/2$である。従って
\begin{equation}
 \Delta T=\frac{\langle P\rangle(1\,\mathrm s)}{C}
 =\frac{\pi^3\mu_0^2H_0^2f^2LR^4}{4\rho C}(1\,\mathrm s).
\end{equation}
与えられた値を代入すると$\langle P\rangle\simeq1.96\times10^{-3}\,\mathrm W$、$\Delta T\simeq3.92\times10^{-3}\,\mathrm K$である。一桁の有効数字では
\begin{equation}
 \boxed{\Delta T\simeq4\times10^{-3}\,\mathrm K}.
\end{equation}

\section*{[2.1] 一つの円環電流の軸上磁場}
半径$a$のループの各電流要素についてBiot--Savart則を用いると、対称性により横成分は相殺し、軸成分だけが残る。距離$\sqrt{a^2+x^2}$、軸への射影因子$a/\sqrt{a^2+x^2}$を使えば
\begin{equation}
 B_x(x)=\frac{\mu_0I}{4\pi}\frac{2\pi a^2}{(a^2+x^2)^{3/2}}
 =\boxed{\frac{\mu_0Ia^2}{2(a^2+x^2)^{3/2}}}.
\end{equation}
ここでは電流方向を右ねじの法則で$+x$向きに磁場を作る向きとした。逆向き電流なら符号が反転する。

\section*{[2.2] 逆向き電流の二つのループ}
左のループを$x=-b$で$+x$向きの磁場を作るように取り、右のループ$x=b$の電流は反対向きとする。軸上の符号付き成分は
\begin{equation}
 B_x(x)=\frac{\mu_0Ia^2}{2}\left[
 \frac{1}{(a^2+(b+x)^2)^{3/2}}-
 \frac{1}{(a^2+(b-x)^2)^{3/2}}\right].
\end{equation}
原点で$B_x(0)=0$であり、奇関数なので偶数次の項は消える。微分すると
\begin{equation}
 B_x'(0)=-\frac{3\mu_0Ia^2b}{(a^2+b^2)^{5/2}},\qquad
 \boxed{B_x(x)=-\frac{3\mu_0Ia^2b}{(a^2+b^2)^{5/2}}x+O(x^3)}.
\end{equation}

\section*{[2.3] 同じ向きの電流と二次項の相殺}
右の電流を反転すると二つの軸成分が加わる。
\begin{equation}
 B_x(x)=\frac{\mu_0Ia^2}{2}\left[
 \frac{1}{(a^2+(b+x)^2)^{3/2}}+
 \frac{1}{(a^2+(b-x)^2)^{3/2}}\right].
\end{equation}
$S=a^2+b^2$とおいて
\begin{equation}
 \frac{1}{[a^2+(b\pm x)^2]^{3/2}}
 =S^{-3/2}\left[1\mp\frac{3b}{S}x
 +\frac{3(4b^2-a^2)}{2S^2}x^2+O(x^3)\right]
\end{equation}
と展開し、$+x$と$-x$の式を足すと
\begin{equation}
 B_x(x)=\frac{\mu_0Ia^2}{S^{3/2}}
 +\frac{3\mu_0Ia^2(4b^2-a^2)}{2S^{7/2}}x^2+O(x^4).
\end{equation}
二次項が消える条件は$4b^2-a^2=0$である。$a,b>0$より
\begin{equation}
 \boxed{a=2b}.
\end{equation}
\end{document}
